En un nucli atòmic radioactiu s'esdevenen dues desintegracions radioactives successives, representades en la gràfica de la figura. En l'eix de les abscisses s'indica el nombre de protons ($Z$) i en l'eix de les ordenades, el nombre de neutrons ($N$) dels elements químics que intervenen en el procés.

\figura[0.45]{fig1.png}

\begin{apartat}{1}
Escriviu les equacions de les dues desintegracions radioactives que es produeixen i digueu com s'anomena cadascuna. Indiqueu el nom, el nombre atòmic i el nombre màssic de tots els elements i de totes les partícules que hi intervenen.
\end{apartat}

\begin{apartat}{1}
Si inicialment tenim $N$ nuclis del primer element i el seu període de semidesintegració és de 10,64 hores, calculeu el temps que haurà de passar perquè es desintegrin un $10,0\,\%$ dels nuclis.
\end{apartat}

\dades{Nombres atòmics d'alguns elements químics: or (Au), 79; mercuri (Hg), 80; tal·li (Tl), 81; plom (Pb), 82; bismut (Bi), 83; poloni (Po), 84; àstat (At), 85.}
